%! The Victory of Orthogonality — Unit 7, Lesson 7.4 — Group Activity (Answer Key)
\documentclass[10pt]{article}
\usepackage{linalg-article}
\usepackage{linalg-key}   % pulls in -boxes; do NOT also load -boxes
\newcommand{\TallMath}[1]{$\displaystyle #1\rule[-1.4em]{0pt}{3.2em}$}
\newcommand{\uu}{\mathbf{u}}
\newcommand{\vv}{\mathbf{v}}
\newcommand{\xx}{\mathbf{x}}
\newcommand{\yy}{\mathbf{y}}
\newcommand{\qq}{\mathbf{q}}
\newcommand{\zero}{\mathbf{0}}
\newcommand{\T}{^{\top}}

\begin{document}

\pageheader{Unit 7, Lesson 7.4}{Group Activity --- Answer Key}
\namepartnerperiod

\vspace{0.10in}

\begin{headlinebox}{blush}
\small\textbf{Orthogonal matrices.} A matrix $Q$ is \textbf{orthogonal} when its columns are perpendicular unit
vectors --- the one test is $Q\T Q=I$. Then the inverse is free ($Q^{-1}=Q\T$) and $Q$ preserves length
($|Q\xx|=|\xx|$). In the SVD $A=U\Sigma V\T$ the orthogonal $U,V$ only rotate; $\Sigma$ does all the stretching.
Work with your partner, show every step, and \emph{justify} each answer.
\end{headlinebox}

\vspace{0.08in}

\begin{tcolorbox}[colback=white, colframe=black!40, title=\textbf{Tier R --- Remediate},
    fonttitle=\bfseries, arc=1.5mm, left=3mm, right=3mm, top=2mm, bottom=2mm, breakable]
Verify that $Q=\tfrac15\begin{bsmallmatrix} 3 & -4\\ 4 & 3 \end{bsmallmatrix}$ is orthogonal.
\begin{enumerate}[leftmargin=1.7em, itemsep=9pt]
  \item \textbf{Unit columns?} The columns are $\qq_1=\tfrac15(3,4)$ and $\qq_2=\tfrac15(-4,3)$. Compute $\qq_1\cdot\qq_1$ and $\qq_2\cdot\qq_2$ --- are both columns length $1$?
        \ansline{$\qq_1\cdot\qq_1=\tfrac1{25}(9+16)=1$ and $\qq_2\cdot\qq_2=\tfrac1{25}(16+9)=1$: yes, both unit length.}
  \item \textbf{Perpendicular?} Compute $\qq_1\cdot\qq_2$.
        \ansline{$\qq_1\cdot\qq_2=\tfrac1{25}(3\cdot(-4)+4\cdot3)=\tfrac1{25}(-12+12)=0$: perpendicular.}
  \item \textbf{The test.} Put those three dot products together as $Q\T Q=\ans{$\begin{bsmallmatrix} 1 & 0\\ 0 & 1 \end{bsmallmatrix}$}$. Is $Q$ orthogonal?
        \ansline{$Q\T Q=I$, so \textbf{yes}, $Q$ is orthogonal.}
\end{enumerate}
\end{tcolorbox}

\begin{tcolorbox}[colback=white, colframe=black!40, title=\textbf{Tier A --- Approaching Proficiency},
    fonttitle=\bfseries, arc=1.5mm, left=3mm, right=3mm, top=2mm, bottom=2mm, breakable]
Keep working with the orthogonal $Q=\tfrac15\begin{bsmallmatrix} 3 & -4\\ 4 & 3 \end{bsmallmatrix}$.
\begin{enumerate}[leftmargin=1.7em, itemsep=9pt]
  \item \textbf{Length preserved.} Let $\xx=(5,0)$, so $|\xx|=5$. Compute $Q\xx$ and its length. Did $Q$ change the length?
        \ansline{$Q\xx=\tfrac15(3\cdot5-4\cdot0,\ 4\cdot5+3\cdot0)=\tfrac15(15,20)=(3,4)$, length $\sqrt{9+16}=5$. \emph{No} --- same length, as an orthogonal matrix must.}
  \item \textbf{Free inverse.} Write $Q^{-1}$ with no computation.
        \ansline{$Q^{-1}=Q\T=\tfrac15\begin{bsmallmatrix} 3 & 4\\ -4 & 3 \end{bsmallmatrix}$.}
  \item \textbf{Interpret (justify).} A rotation has $\det=+1$. Compute $\det Q$, and explain in one sentence what $Q$ does to space geometrically.
        \ansline{$\det Q=\tfrac1{25}(3\cdot3-(-4)\cdot4)=\tfrac{25}{25}=1$: $Q$ is a \textbf{rotation} --- it turns every vector by the same angle ($\cos\theta=\tfrac35$) without changing any length.}
\end{enumerate}
\end{tcolorbox}

\begin{tcolorbox}[colback=white, colframe=black!40, title=\textbf{Tier E --- Extension},
    fonttitle=\bfseries, arc=1.5mm, left=3mm, right=3mm, top=2mm, bottom=2mm, breakable]
\textbf{Rotate--stretch--rotate.} A matrix $A=U\Sigma V\T$ has orthogonal $U,V$ and
$\Sigma=\begin{bsmallmatrix} 4 & 0\\ 0 & 1 \end{bsmallmatrix}$ (so its singular values are $\sigma_1=4,\ \sigma_2=1$).
\begin{enumerate}[leftmargin=1.7em, itemsep=9pt]
  \item \textbf{Biggest stretch.} As $\xx$ runs over all unit vectors, what is the \emph{largest} value of $|A\xx|$? Along which singular direction is it reached?
        \ansline{$|A\xx|$ is largest $=\sigma_1=4$, reached along the first input singular direction $\vv_1$ (there $A\vv_1=\sigma_1\uu_1$, length $4$).}
  \item \textbf{Smallest stretch.} What is the \emph{smallest} value of $|A\xx|$ for a unit vector $\xx$?
        \ansline{The smallest is $\sigma_2=1$, reached along $\vv_2$.}
  \item \textbf{Why (justify).} Explain why $U$ and $V\T$ don't change the stretch factor --- so all of the stretching comes from $\Sigma$.
        \ansline{$U$ and $V\T$ are orthogonal, so they preserve length ($|V\T\xx|=|\xx|$ and $|U\yy|=|\yy|$). Only $\Sigma$ scales lengths --- by $\sigma_1$ and $\sigma_2$ along the axes --- so the entire stretch of $A$ is set by the singular values, between $\sigma_2$ (smallest) and $\sigma_1$ (largest).}
\end{enumerate}
\end{tcolorbox}

\begin{teachernote}
Tier~R walks the definition one dot product at a time on the clean $3$--$4$--$5$ rotation
$Q=\tfrac15\begin{bsmallmatrix} 3 & -4\\ 4 & 3 \end{bsmallmatrix}$: unit columns ($\tfrac1{25}(9+16)=1$),
perpendicular ($-12+12=0$), hence $Q\T Q=I$. Tier~A uses the \emph{same} $Q$ to hit the two payoffs --- length
preservation ($\xx=(5,0)\to Q\xx=(3,4)$, still length $5$) and the free inverse $Q^{-1}=Q\T$ --- then reads
$\det Q=1$ as ``rotation.'' Tier~E is the conceptual capstone: with orthogonal $U,V$, the stretch factor of
$A=U\Sigma V\T$ runs from $\sigma_2=1$ up to $\sigma_1=4$ because only $\Sigma$ scales lengths --- this is the
``all the stretching lives in $\Sigma$'' idea made precise. If time is short, Tiers~R and~A are the must-dos.
Watch for: checking perpendicularity but forgetting the unit-length check; recomputing $Q^{-1}$ instead of just
transposing; thinking $U$ or $V$ could add stretch.
\end{teachernote}

\end{document}
